\chapter{Project Statement and Scope}
\label{ch:statement}

\section{The project as originally stated}

The project was defined as follows.

\begin{quote}
\textbf{Project Title:} Solar Panel Maximum Power Point Tracking as a Control
Problem.

Core idea is to use feedback control to maintain operations near maximum power
despite changing irradiance. In MATLAB, need to model PV model, MPPT model, and
control analysis. In physical model, need solar panel and DC-DC model.
\end{quote}

Everything in this report follows from that statement. This first chapter
unpacks it: what each phrase means, what a complete answer has to contain, and
how the work was organised to deliver it.

\section{Reading the statement}

\subsection{``Maximum power point''}

A solar panel is neither a voltage source nor a current source. Its terminal
current depends nonlinearly on its terminal voltage, and the product of the two,
the delivered power, has a single interior maximum. That maximum is the
\term{maximum power point} (MPP). Operating anywhere else wastes energy that the
sunlight already delivered to the panel.

\subsection{``Despite changing irradiance''}

The location of the MPP is not fixed. It moves when the \term{irradiance}, the
optical power per unit area reaching the panel, changes because of clouds, time
of day or shading. It also moves when the cell temperature changes. So the
problem is not to find a number once; it is to \emph{track} a moving target
continuously, in real time, using only measurements of terminal voltage and
terminal current.

\subsection{``As a control problem''}

This is the phrase that sets the character of the whole project. Much of the
published work on MPPT treats it as an \emph{algorithm} question: which search
rule finds the peak fastest. That framing hides the fact that the search rule
sits on top of a physical plant, a switching power converter, whose dynamics
decide whether the search rule can work at all.

Treating MPPT as a control problem means asking the questions a control engineer
would ask:

\begin{itemize}[leftmargin=1.4em]
  \item What exactly is the plant, and what is its transfer function?
  \item How do the plant dynamics change as the operating point moves along the
        I--V curve and as irradiance changes?
  \item What is the control input, what is the measured output, and what is the
        error signal?
  \item What stability margins does the closed loop have, and are they
        maintained over the whole operating envelope?
  \item What fundamental limits, right half plane zeros, resonances, sampling
        delay, bound the achievable bandwidth?
  \item Given those limits, how fast can a tracker legitimately be, and what
        does it cost in steady-state accuracy?
\end{itemize}

A central result of this project, developed in
Chapter~\ref{ch:converter} and Chapter~\ref{ch:control}, is that these questions
have sharp and somewhat surprising answers, and that the answers explain the
behaviour of the classical MPPT algorithms rather than merely sitting alongside
them.

\subsection{``In MATLAB, need to model PV model, MPPT model, and control analysis''}

Three deliverables on the analysis side:

\begin{enumerate}[leftmargin=1.6em]
  \item A \term{photovoltaic model}: a mathematical description of the panel
        accurate enough that its curvature, not just its three datasheet points,
        can be trusted. Chapter~\ref{ch:pvmodel}.
  \item An \term{MPPT model}: implementations of the tracking algorithms as
        proper discrete-time controllers with defined sample rates and states,
        so that they can be compared fairly. Chapter~\ref{ch:mppt}.
  \item A \term{control analysis}: linearisation, frequency response, stability
        margins and robustness across the operating envelope.
        Chapters~\ref{ch:converter} and~\ref{ch:control}.
\end{enumerate}

\subsection{``In physical model, need solar panel and DC-DC model''}

A separate deliverable: a model of the actual hardware, not an idealised
abstraction. The panel appears as a physical two-terminal electrical device, and
the \term{DC--DC converter}, the switching circuit that transforms one direct
voltage into another, appears as a real circuit with a transistor that switches
at \SI{20}{\kilo\hertz}, a diode, an inductor with resistance and capacitors.
Chapter~\ref{ch:physical} covers this, together with the hardware required to
build it.

\section{What a complete answer requires}

Putting those readings together, a complete project needs all of the following,
and each is delivered in the chapter named.

\begin{longtable}{@{}L{0.40\textwidth}L{0.42\textwidth}L{0.12\textwidth}@{}}
\toprule
\textbf{Requirement} & \textbf{How it is met} & \textbf{Chapter} \\
\midrule
\endfirsthead
\toprule
\textbf{Requirement} & \textbf{How it is met} & \textbf{Chapter} \\
\midrule
\endhead
\bottomrule
\endfoot
A panel model valid away from the datasheet points & Five-parameter single-diode model with De Soto translation, parameters identified from five simultaneous datasheet constraints & \ref{ch:pvmodel} \\
Ground truth to score trackers against & Exact maximum power point by root solve on the analytic power gradient & \ref{ch:pvmodel} \\
A converter model suitable for control design & State-space averaged boost converter, linearised analytically, Jacobian verified against finite differences & \ref{ch:converter} \\
Understanding of how the plant changes & Operating envelope sweep in irradiance and along the I--V curve & \ref{ch:converter} \\
A controller with quantified margins & Cascaded current and voltage loops, margins computed over the envelope, compared against a single-loop baseline & \ref{ch:control} \\
MPPT algorithms implemented fairly & Common tracker interface, identical sample rate and identical inner control loops & \ref{ch:mppt} \\
An objective comparison metric & Dynamic MPPT efficiency on EN 50530-style irradiance ramps & \ref{ch:benchmark} \\
Reproducibility and correctness evidence & Scripted workflow and 27 automated unit tests & \ref{ch:software} \\
A physical circuit model & Switched Simscape model with a custom photovoltaic component, cross-validated against the averaged model & \ref{ch:physical} \\
A route to hardware & Component selection, ratings, sensing, protection, and a staged build plan & \ref{ch:physical} \\
\end{longtable}

\section{System chosen for study}

All results in this report are for one concrete system, so that numbers can be
compared directly from chapter to chapter.

\begin{table}[htbp]
\centering
\caption{The system under study.}
\label{tab:system}
\begin{tabular}{@{}lll@{}}
\toprule
\textbf{Item} & \textbf{Value} & \textbf{Note} \\
\midrule
Module & Kyocera KC200GT & \SI{200}{\watt}, 54 cells, multicrystalline silicon \\
Array & 6 series $\times$ 2 parallel & \SI{2402}{\watt} at STC \\
Array MPP at STC & \SI{157.8}{\volt}, \SI{15.22}{\ampere} & \\
Converter & Boost & PV side to DC bus \\
Inductance $L$ & \SI{1.5}{\milli\henry} & \SI{20}{\percent} current ripple \\
Input capacitance $C_{\mathrm{pv}}$ & \SI{100}{\micro\farad} & sets PV voltage ripple \\
Inductor resistance $R_L$ & \SI{0.15}{\ohm} & \\
Switching frequency $f_{sw}$ & \SI{20}{\kilo\hertz} & \\
DC bus $V_{dc}$ & \SI{400}{\volt} & treated as stiff \\
\bottomrule
\end{tabular}
\end{table}

The KC200GT was chosen because its datasheet is widely used as a modelling
benchmark in the literature~\cite{villalva}, which makes the extracted
parameters checkable against published values. The array size, roughly
\SI{2.4}{\kilo\watt} into a \SI{400}{\volt} bus, is representative of the
photovoltaic input stage of a small single-phase string inverter.

\section{The control architecture adopted}

The architecture that emerges from the analysis, and which the rest of the
report develops and justifies, is a three-layer hierarchy separated by
timescale.

\begin{table}[htbp]
\centering
\caption{The three-layer control hierarchy and its timescale separation.}
\label{tab:hierarchy}
\begin{tabular}{@{}llll@{}}
\toprule
\textbf{Layer} & \textbf{Controlled variable} & \textbf{Command it issues} & \textbf{Bandwidth} \\
\midrule
MPPT & PV power (indirectly) & voltage reference $\Vmp^\ast$ & \SI{50}{\hertz} \\
Outer loop & PV voltage $\vpv$ & inductor current reference & \SI{250}{\hertz} \\
Inner loop & inductor current $i_L$ & duty ratio $d$ & \SI{2}{\kilo\hertz} \\
\bottomrule
\end{tabular}
\end{table}

Each layer is roughly eight times slower than the one beneath it. That
separation is what allows each to be designed as though the layer below were
instantaneous and the layer above were frozen, and Chapter~\ref{ch:control}
verifies that the assumption holds by computing margins on the full-order model
rather than on the simplified one used for design.
